% This is continuation file - append to main document

\subsubsection{Strategy 3: Simple Chunking (Fallback)}

If LLM inference fails, fall back to simple page-based chunking:

\begin{algorithm}[H]
\caption{Simple Page Chunking}
\begin{algorithmic}[1]
\Require Pages $P$, Chunk size $m$
\Ensure Tree root $r$

\State $r \gets \text{CreateNode}(\text{``Document''}, 1, |P|, 0)$

\For{$i = 1$ to $|P|$ step $m$}
    \State $p_{\text{start}} \gets i$
    \State $p_{\text{end}} \gets \min(i + m - 1, |P|)$
    \State $\text{content} \gets \text{ConcatenatePages}(P, p_{\text{start}}, p_{\text{end}})$
    \State $v \gets \text{CreateNode}(\text{``Section ''} || (i/m + 1), p_{\text{start}}, p_{\text{end}}, 1, \text{content})$
    \State $r.\text{children.append}(v)$
\EndFor

\State \Return $r$
\end{algorithmic}
\end{algorithm}

\subsection{Summary Generation}

\subsubsection{Node-Level Summarization}

For each node, optionally generate an LLM summary to aid retrieval:

\begin{algorithm}[H]
\caption{Add Summaries to Tree}
\begin{algorithmic}[1]
\Require Tree node $v$
\Ensure Node $v$ with summary

\If{$v.\text{content} \neq \emptyset$ \textbf{and} $v.\text{summary} = \emptyset$}
    \State $\text{prompt} \gets \text{BuildSummaryPrompt}(v.\text{title}, v.\text{content})$
    \State $v.\text{summary} \gets \text{LLM}(\text{prompt})$
\EndIf

\For{each child $c$ in $v.\text{children}$}
    \State $\text{AddSummaries}(c)$
\EndFor
\end{algorithmic}
\end{algorithm}

\begin{conceptbox}{Why Summaries?}
Node summaries serve two purposes:
\begin{enumerate}
    \item \textbf{Improved Retrieval}: LLM can match query to summary without reading full content
    \item \textbf{User Preview}: Users see what each section contains before diving into details
\end{enumerate}

Trade-off: Summaries increase indexing time but improve retrieval quality.
\end{conceptbox}

% ============================================================================
\section{LLM-Based Retrieval}
% ============================================================================

\subsection{Tree Representation for Search}

\subsubsection{Converting Tree to Text}

The tree must be converted to a text representation the LLM can reason over:

\begin{algorithm}[H]
\caption{Tree to Search Representation}
\begin{algorithmic}[1]
\Require Tree node $v$, Depth $d$
\Ensure Text representation $\text{repr}$

\State $\text{indent} \gets \text{``  ''} \times d$
\State $\text{repr} \gets \text{indent} || \text{``[''} || v.\text{id} || \text{``] ''} || v.\text{title}$
\State $\text{repr} \gets \text{repr} || \text{`` (pages ''} || v.p_{\text{start}} || \text{``-''} || v.p_{\text{end}} || \text{``)''}$

\If{$v.\text{summary} \neq \emptyset$}
    \State $\text{repr} \gets \text{repr} || \text{``\textbackslash n''} || \text{indent} || \text{``  Summary: ''} || v.\text{summary}$
\EndIf

\For{each child $c$ in $v.\text{children}$}
    \State $\text{repr} \gets \text{repr} || \text{``\textbackslash n''} || \text{TreeToText}(c, d+1)$
\EndFor

\State \Return $\text{repr}$
\end{algorithmic}
\end{algorithm}

\textbf{Example Output}:
\begin{lstlisting}
[node_0001] Document (pages 1-50)
  [node_0002] Introduction (pages 1-5)
    Summary: Overview of research objectives and methodology...
  [node_0003] Literature Review (pages 6-12)
    Summary: Summary of related work in machine learning...
  [node_0004] Methodology (pages 13-25)
    [node_0005] Data Collection (pages 13-17)
      Summary: Description of dataset sources...
\end{lstlisting}

\subsection{Search Prompt Construction}

\subsubsection{Base Search Prompt}

The search prompt instructs the LLM to reason about relevance:

\begin{lstlisting}[language=Python, caption=Search Prompt Template]
base_prompt = f"""You are a document search expert.
Given a question and document tree structure,
identify the most relevant sections.

QUESTION:
{query}

DOCUMENT STRUCTURE:
{tree_repr}

TASK:
1. Reason about which sections likely contain
   information to answer the question
2. Consider section titles, summaries, and
   hierarchical relationships
3. Return up to {max_nodes} most relevant node IDs

Respond ONLY with JSON in this format:
{{
  "reasoning": "Your step-by-step reasoning...",
  "node_ids": ["node_0001", "node_0002", ...]
}}
"""
\end{lstlisting}

\subsubsection{Domain Rules Integration}

Expert knowledge can be injected via additional instructions:

\begin{lstlisting}[language=Python, caption=Domain Rules Example]
domain_rules = """
EXPERT RULES FOR FINANCIAL DOCUMENTS:
1. For revenue/income questions, prioritize:
   - Income Statement
   - MD&A (Management Discussion & Analysis)
2. For risk questions, check:
   - Risk Factors section
   - MD&A
3. For balance sheet items, go to:
   - Balance Sheet
   - Notes to Financial Statements
"""

prompt = base_prompt + domain_rules
\end{lstlisting}

\subsection{Retrieval Examples}

\begin{examplebox}{Example 1: Direct Question}
\textbf{Query}: ``What was the revenue in Q3 2023?''

\textbf{LLM Reasoning}:
\begin{itemize}
    \item Revenue is financial metric
    \item Q3 2023 is specific time period
    \item Should check: Income Statement, MD\&A
\end{itemize}

\textbf{Retrieved Nodes}:
\begin{itemize}
    \item [node\_0015] Income Statement (pages 22-24)
    \item [node\_0018] Q3 2023 Financial Highlights (pages 8-10)
\end{itemize}
\end{examplebox}

\begin{examplebox}{Example 2: Multi-Hop Question}
\textbf{Query}: ``How did changes in foreign exchange rates affect operating expenses?''

\textbf{LLM Reasoning}:
\begin{itemize}
    \item Requires: FX rate changes + operating expenses
    \item FX discussion likely in: MD\&A, Risk Factors
    \item Operating expenses in: Income Statement, MD\&A
    \item Need both sections for complete answer
\end{itemize}

\textbf{Retrieved Nodes}:
\begin{itemize}
    \item [node\_0012] MD\&A: Foreign Exchange Impact (pages 15-17)
    \item [node\_0014] Operating Expenses Analysis (pages 19-21)
    \item [node\_0025] Risk Factors: Currency Risk (pages 38-40)
\end{itemize}
\end{examplebox}

% ============================================================================
\section{Answer Generation}
% ============================================================================

\subsection{Context Assembly}

After retrieval, assemble context from retrieved nodes:

\begin{algorithm}[H]
\caption{Assemble Context}
\begin{algorithmic}[1]
\Require Retrieved nodes $N_{\text{ret}}$, Tree $\mathcal{T}$
\Ensure Context string $\text{ctx}$

\State $\text{ctx} \gets \text{``''}$

\For{each node ID $\text{nid}$ in $N_{\text{ret}}$}
    \State $v \gets \text{FindNode}(\mathcal{T}, \text{nid})$
    
    \State $\text{section} \gets \text{``\textbackslash nSection: ''} || v.\text{title}$
    \State $\text{section} \gets \text{section} || \text{``\textbackslash nPages: ''} || v.p_{\text{start}} || \text{``-''} || v.p_{\text{end}}$
    \State $\text{section} \gets \text{section} || \text{``\textbackslash nContent:\textbackslash n''} || v.\text{content}$
    
    \State $\text{ctx} \gets \text{ctx} || \text{section} || \text{``\textbackslash n---\textbackslash n''}$
\EndFor

\State \Return $\text{ctx}$
\end{algorithmic}
\end{algorithm}

\subsection{Citation-Aware Generation}

\subsubsection{Answer Generation Prompt}

\begin{lstlisting}[language=Python, caption=Answer Generation Template]
prompt = f"""Answer the question using ONLY
the provided context. Cite specific sections
in your answer.

CONTEXT:
{context}

QUESTION:
{question}

INSTRUCTIONS:
1. Answer clearly and concisely
2. Cite sections by title and page numbers
3. If context doesn't contain enough
   information, say so
4. Use direct quotes when relevant

ANSWER:"""
\end{lstlisting}

\subsubsection{Citation Format}

The generated answer should include citations:

\begin{examplebox}{Example Answer with Citations}
\textbf{Question}: ``What was Q3 2023 revenue?''

\textbf{Answer}:
Q3 2023 revenue was \$42.3 million (\textit{Income Statement}, pages 22-24), representing a 15\% year-over-year increase. This growth was primarily driven by expansion in the Asia-Pacific market (\textit{MD\&A: Revenue Analysis}, pages 10-12).
\end{examplebox}

% ============================================================================
\section{Supervised Learning \& Evaluation}
% ============================================================================

\subsection{Creating Supervised Dataset}

\subsubsection{QA Pair Structure}

Each supervised example consists of:

\begin{lstlisting}[language=json, caption=Supervised QA Pair Format]
{
  "question": "What were the main risk factors?",
  "ground_truth_answer": "The main risks were: ...",
  "relevant_node_ids": ["node_0025", "node_0026"],
  "context_needed": "Risk Factors section",
  "difficulty": "medium",
  "category": "risk_analysis"
}
\end{lstlisting}

\subsubsection{Dataset Creation Process}

\begin{algorithm}[H]
\caption{Create Supervised Dataset}
\begin{algorithmic}[1]
\Require Document $D$, Tree $\mathcal{T}$
\Ensure QA pairs $\mathcal{Q}$

\State $\mathcal{Q} \gets []$

\Comment{Manual annotation}
\For{each important section $s$ in $D$}
    \State Formulate question $q$ that section answers
    \State Extract ground truth answer $a^*$ from section
    \State Identify node IDs $N^*$ containing answer
    \State Assign difficulty level $\delta \in \{\text{easy, medium, hard}\}$
    \State Create QA pair $\langle q, a^*, N^*, \delta \rangle$
    \State $\mathcal{Q}.\text{append}(\text{qa})$
\EndFor

\State \Return $\mathcal{Q}$
\end{algorithmic}
\end{algorithm}

\subsection{Evaluation Protocol}

\subsubsection{Full Evaluation Pipeline}

\begin{algorithm}[H]
\caption{Evaluate RAG System}
\begin{algorithmic}[1]
\Require RAG system $S$, QA pairs $\mathcal{Q}$
\Ensure Evaluation metrics $M$

\State $\text{results} \gets []$

\For{each $\langle q, a^*, N^*, \delta \rangle$ in $\mathcal{Q}$}
    \Comment{Run RAG query}
    \State $N_{\text{ret}}, \hat{a} \gets S.\text{query}(q)$
    
    \Comment{Evaluate retrieval}
    \State $m_{\text{ret}} \gets \text{EvaluateRetrieval}(N_{\text{ret}}, N^*)$
    
    \Comment{Evaluate answer}
    \State $m_{\text{ans}} \gets \text{EvaluateAnswer}(\hat{a}, a^*)$
    
    \State $\text{results.append}(m_{\text{ret}}, m_{\text{ans}})$
\EndFor

\Comment{Aggregate metrics}
\State $M \gets \text{AggregateMetrics}(\text{results})$

\State \Return $M$
\end{algorithmic}
\end{algorithm}

\subsection{Failure Analysis}

\subsubsection{Identifying Patterns}

When evaluation reveals low performance, analyze failure patterns:

\begin{algorithm}[H]
\caption{Analyze Failures}
\begin{algorithmic}[1]
\Require Evaluation results $R$, Threshold $\theta$
\Ensure Failure patterns $F$

\State $F \gets []$

\For{each result $r$ in $R$}
    \If{$r.F_1 < \theta$}
        \State Extract: $q$, $N_{\text{ret}}$, $N^*$, difficulty
        \State $F.\text{append}(\text{failure case})$
    \EndIf
\EndFor

\Comment{Use LLM to identify patterns}
\State $\text{prompt} \gets \text{BuildAnalysisPrompt}(F)$
\State $\text{analysis} \gets \text{LLM}(\text{prompt})$

\State \Return $\text{analysis}$
\end{algorithmic}
\end{algorithm}

\subsubsection{Domain Rule Suggestions}

The LLM analyzes failures and suggests domain rules:

\begin{examplebox}{Failure Analysis Example}
\textbf{Pattern Identified}: 
System consistently misses nodes in ``Notes to Financial Statements'' when answering accounting policy questions.

\textbf{Suggested Rule}:
``For questions about accounting policies, depreciation methods, or revenue recognition, always check the Notes to Financial Statements section in addition to the main financial statements.''
\end{examplebox}

% ============================================================================
\section{Implementation Details}
% ============================================================================

\subsection{System Architecture}

\begin{figure}[H]
\centering
\begin{tikzpicture}[
    node distance=1.5cm,
    box/.style={rectangle, draw, fill=blue!20, text width=3cm, align=center, rounded corners},
    arrow/.style={->, >=stealth, thick}
]

% Nodes
\node[box] (pdf) {PDF Document};
\node[box, below of=pdf] (extract) {PDF Processor};
\node[box, below of=extract] (tree) {Tree Builder};
\node[box, below of=tree] (index) {Tree Index};

\node[box, right=4cm of index] (query) {User Query};
\node[box, above of=query] (search) {LLM Tree Searcher};
\node[box, above of=search] (gen) {Answer Generator};
\node[box, above of=gen] (answer) {Final Answer};

\node[box, right=4cm of tree] (sup) {Supervised Dataset};
\node[box, below of=sup] (eval) {Evaluator};

% Arrows
\draw[arrow] (pdf) -- (extract);
\draw[arrow] (extract) -- (tree);
\draw[arrow] (tree) -- (index);

\draw[arrow] (query) -- (search);
\draw[arrow] (index) -- (search);
\draw[arrow] (search) -- node[right] {Retrieved Nodes} (gen);
\draw[arrow] (index) -- (gen);
\draw[arrow] (gen) -- (answer);

\draw[arrow] (sup) -- (eval);
\draw[arrow] (index) -- (eval);
\draw[arrow] (eval) -- node[below] {Metrics \& Feedback} (tree);

\end{tikzpicture}
\caption{Vectorless RAG System Architecture}
\end{figure}

\subsection{Key Classes}

\subsubsection{TreeNode Class}

\begin{lstlisting}[language=Python, caption=TreeNode Implementation]
@dataclass
class TreeNode:
    """Node in document tree"""
    node_id: str
    title: str
    page_start: int
    page_end: int
    level: int
    summary: Optional[str] = None
    content: Optional[str] = None
    children: List['TreeNode'] = None
    
    def __post_init__(self):
        if self.children is None:
            self.children = []
    
    def to_dict(self) -> Dict:
        """Serialize to JSON"""
        return {
            'node_id': self.node_id,
            'title': self.title,
            'page_start': self.page_start,
            'page_end': self.page_end,
            'level': self.level,
            'summary': self.summary,
            'content': self.content,
            'children': [
                child.to_dict()
                for child in self.children
            ]
        }
\end{lstlisting}

\subsubsection{VectorlessRAG Class}

\begin{lstlisting}[language=Python, caption=Main RAG Class]
class VectorlessRAG:
    """Main RAG system"""
    
    def __init__(self, openai_api_key: str):
        self.client = OpenAI(api_key=openai_api_key)
        self.pdf_processor = PDFProcessor()
        self.tree_builder = TreeBuilder(self.client)
        self.searcher = LLMTreeSearcher(self.client)
        self.tree = None
    
    def index_document(
        self,
        pdf_path: str,
        cache_path: Optional[str] = None
    ) -> TreeNode:
        """Index PDF and build tree"""
        # Extract pages
        pages = self.pdf_processor.extract_pages(
            pdf_path
        )
        
        # Build tree
        self.tree = self.tree_builder.build_tree(
            pages
        )
        
        # Cache if requested
        if cache_path:
            self._save_tree(cache_path)
        
        return self.tree
    
    def query(
        self,
        question: str,
        max_nodes: int = 5
    ) -> Dict[str, Any]:
        """Query indexed document"""
        # Search tree
        node_ids = self.searcher.search(
            question, self.tree, max_nodes
        )
        
        # Retrieve nodes
        nodes = self._retrieve_nodes(node_ids)
        
        # Generate answer
        answer = self._generate_answer(
            question, nodes
        )
        
        return {
            'answer': answer,
            'sources': [
                {
                    'title': n.title,
                    'pages': f"{n.page_start}-{n.page_end}"
                }
                for n in nodes
            ],
            'node_ids': node_ids
        }
\end{lstlisting}

\subsection{Configuration}

\subsubsection{Environment Variables}

\begin{lstlisting}[language=bash, caption=.env File]
# OpenAI API Key
OPENAI_API_KEY=sk-your-key-here

# Model selection
LLM_MODEL=gpt-4o-mini

# Tree building parameters
MAX_PAGES_PER_NODE=10
TOC_CHECK_PAGES=20
GENERATE_SUMMARIES=true

# Retrieval parameters
MAX_RETRIEVED_NODES=5
\end{lstlisting}

% ============================================================================
\section{Usage Examples}
% ============================================================================

\subsection{Basic Usage}

\subsubsection{Complete Example}

\begin{lstlisting}[language=Python, caption=Basic Usage Example]
from vectorless_rag import VectorlessRAG
import os

# Initialize system
rag = VectorlessRAG(
    openai_api_key=os.getenv("OPENAI_API_KEY")
)

# Index a document
tree = rag.index_document(
    pdf_path="research_paper.pdf",
    cache_path="research_paper_tree.json"
)

# Print tree structure
rag.print_tree()

# Query the document
result = rag.query(
    "What methodology was used?",
    max_nodes=3
)

print("Answer:", result['answer'])
print("\nSources:")
for source in result['sources']:
    print(f"  - {source['title']} (pp. {source['pages']})")
\end{lstlisting}

\subsection{Supervised Evaluation}

\begin{lstlisting}[language=Python, caption=Evaluation Example]
from supervised_rag import (
    SupervisedRAGEvaluator,
    SupervisedQAPair
)

# Create evaluator
evaluator = SupervisedRAGEvaluator(
    openai_client=rag.client
)

# Add QA pairs
evaluator.add_qa_pair(
    question="What was the sample size?",
    ground_truth_answer="The study used 500 participants",
    relevant_node_ids=["node_0005"],
    context_needed="Methodology section",
    difficulty="easy",
    category="methods"
)

# Save for reuse
evaluator.save_qa_pairs("qa_dataset.json")

# Run evaluation
results = evaluator.evaluate_rag_system(
    rag_system=rag,
    verbose=True
)

# View metrics
print(results['aggregate_metrics']['summary'])
\end{lstlisting}

\end{document}
